% ===========================================================================
% 04b-representation.tex -- beats 3-4. The expression frame, read in the weights.
% Source: thesis/Sections/Investigation-v4.tex 1052-1545.
% Every number, interval, n and conclusion is verbatim from v1.
%
% Consistency pass: the section now opens by naming the third question of
% \cref{sec:intro-map} and stating its answer -- identity written into the
% activations, barely consulted by generation -- before the instruments, so a
% reader arriving from the introduction finds what was promised. The stratum
% paragraph and its \scopelimit are adjacent again, since the scope limit is
% about that stratum's reference. "Step three, the removal gate" was a \defn
% box between three \paragraph steps; it is now the paragraph it always was,
% merged with the gate result that followed it. Vocabulary: "output" as a noun
% for the prediction, "compound", and bare "signature" are gone; "shift" is
% reserved for treated minus control, so the logsumexp constant is an offset.
% ===========================================================================

% ===========================================================================
\beatsection{E}{3}{The drug is not missing from the computation}{sec:methods-probe}

\begin{question}[3]
A model that ignores a drug it was told about can fail in two ways. Either the
drug never enters its computation, an activation-side encoding failure that
target- and input-side fixes might reach, or it enters and nothing reads it, a
decoder failure only an objective-side fix could reach. Behaviour cannot tell
them apart, since both predict the same scores. The weights can. A linear decode
probe, read as a floor test and not as evidence of use, settles the first half.
\end{question}

\noindent
\Cref{sec:intro-map} asks where useful prediction fails: is the drug absent from
the data, lost from the model's representation, or encoded and underused by
generation? The three sections that follow answer in one direction, and the
answer is worth having before the instruments. Drug identity is written into the
activations and stays written, more legibly with depth once batch structure is
projected out, while generation consults it too weakly to account for the
drug-blindness of \cref{sec:res-blind}. The failure is not on the activation
side. Neither half comes free: the presence result is close to guaranteed by
construction, the low-gain reading rests on the weakest measurement in the
chapter, and both are carried below with the discount attached.

Every behavioural measurement so far fits the last two readings
equally.\framecue{\Eframe}{scores here are full-profile, \cref{sec:res-metrics}}
\Cref{sec:res-blind} has answered the first where it matters here: the task is
empirically discriminable for a substantial stratum at the tested pseudobulk
budget, the reference rises with aggregation, $46.2\%$ of conditions clear the
declared identifiability bar, and on those the model is matched exactly by a
baseline copying the untreated cell ($0.768$ against $0.766$).\gloss{identifiable
stratum}{conditions whose own within-well split-half reference clears the pre-set
bar} The shortfall there is the model's; where in the model it lives needs the
weights.

\begin{scopelimit}[carried into \cref{sec:res-epsilon} and \cref{sec:limitations}]
The distance from a score to that stratum's reference is not available headroom.
Across the three difficulty strata of \cref{sec:res-blind} the control-copy's own
score tracks the reference almost step for step: $0.310$ against a reference of
$0.287$ on the low-reference stratum, $0.540$ against $0.689$ on the marginal one,
and $0.766$ against $0.938$ on the identifiable one. A predictor that reproduces
the untreated cell and knows nothing whatever about any drug does best exactly
where the ranking problem is easiest. What separates the strata on this
criterion is largely drug-independent: within a comparison set every condition's
untreated cells come from one plate-matched DMSO pool, so a control-copy scores
high exactly where a condition's measured profile sits closer to that shared
untreated state than its plate-mates' do. Most of the ranking these conditions
are scored on can therefore be done without the drug, and the gap between a predictor and a high reference is
not all of it room that better drug modelling would fill. Every score quoted on
the identifiable stratum below inherits that reading, including the three target
constructions of \cref{sec:res-epsilon}.
\end{scopelimit}

\paragraph{Three sub-questions, in the order their blind spots close} Is drug
identity present in the activations? Does the network use the region where it is
written? And does it use \emph{which} drug? Each instrument's blind spot is what
the next must close; \cref{sec:res-epsilon} adds a fourth.
\Cref{fig:comparators} maps all four, on scales that may not be read across.

% ===========================================================================
\begin{figure}[htbp]
\begin{fullwidth}
\centering
% ===========================================================================
% figs/comparators.tex -- fig:comparators, reworked for thesis_v2.
% FIGURE BODY ONLY. No preamble, no document environment. \input it from
% inside a figure/fullwidth float in Sections/04b-representation.tex.
%
% Four instruments, four deliberately incommensurable value-to-position maps,
% one comparator each. The premise and all four scales are v1's and unchanged.
% What changed is labelling only; no mark moved, no number changed.
%
%   (1) ROW B plotted seven marks and named one ("all seven layers"), leaving
%       six unidentifiable. Every mark is now named by the layer it is, with a
%       fanned leader for the three that crowd at 3.6/3.8/4.2, and the stub
%       carries "one mark per layer" so a bare accent numeral above the axis
%       cannot be misread as a second set of tick labels (ticks are quietink
%       and below the axis; layer numbers are accent and above it).
%   (2) ROW A's "slab projected out" sat to the right of its own point and
%       nearer the shuffled-label rule than to the mark it names. It now has a
%       leader back to the 0.121 dot.
%   (3) ROW C's single-cell marker was a hollow ring sitting on the dotted zero
%       rule (-0.0012 against 0.000: 0.10cm apart on this scale). It is now
%       solid, the zero rule is interrupted across its height so nothing runs
%       into it, and its label starts 0.30cm clear with a leader through the
%       interruption -- necessary because the zero rule lies between mark and
%       label and would otherwise claim the label.
%   (4) Bar labels in ROWS C and D moved from a 0.13cm to a 0.30cm gap, so
%       "optimal-transport barycentric" no longer reads as a continuation of
%       its own interval bar.
%   General rule applied throughout: a label more than ~0.3cm from its mark
%   gets a short leader.
%   Row heights grew (rows now 0.70/0.77/0.98 apart rather than colliding) to
%   pay for the new labels. Nothing is rescaled: 1cm here is 1cm on the page.
%
% PROVENANCE of every number drawn is in comparators.json, sibling to this
% file. All four rows were checked against RESULTS_cluster and against the
% thesis text; they agree. Sources:
%   row A, B  RESULTS_cluster/workspace_probe.json
%   row C     RESULTS_cluster/probe_arm_residual.json,
%             probe_rep_residual_s43.json, probe_rep_residual_s45.json,
%             probe_arm_single_cell.json
%   row D     RESULTS_cluster/stratify_sameplate_final.json,
%             stratify_ot_T2.json, stratify_consensus.json
% ===========================================================================
\begin{tikzpicture}[x=1cm, y=1cm, font=\small]
  % Four deliberately different value-to-position maps, one per rule.
  \def\xa#1{{(#1)*10.5}}
  \def\xb#1{{(#1)/20*10.5}}
  \def\xc#1{{((#1)+0.02)/0.13*10.5}}
  \def\xd#1{{((#1)+0.03)/0.08*10.5}}
  % Same maps shifted by a label gap; the shift MUST sit inside the braces --
  % \xa{v}+0.05 is not a pgfmath expression and does not compile.
  % 0.30cm, not v1's 0.13cm: at 0.13cm a label abutted the end of its own
  % interval bar and read as part of it.
  \def\xcR#1{{((#1)+0.02)/0.13*10.5+0.30}}
  \def\xdR#1{{((#1)+0.03)/0.08*10.5+0.30}}
  % leader/label shims: +0.05 and +0.06 clear a 1.4pt marker, -0.10 and +0.10
  % keep row A's two abutting labels off their own dots, +0.35 puts row C's
  % single-cell label 0.30cm clear of the marker's edge.
  \def\xaN#1{{(#1)*10.5+0.05}}
  \def\xaL#1{{(#1)*10.5-0.10}}
  \def\xaP#1{{(#1)*10.5+0.10}}
  \def\xcN#1{{((#1)+0.02)/0.13*10.5+0.06}}
  \def\xcM#1{{((#1)+0.02)/0.13*10.5+0.35}}

  % row dividers
  \foreach \y in {-0.70,-2.76,-5.82}{
    \draw[rulegrey!55, line width=0.3pt] (-4.45,\y) -- (11.20,\y);}

  % ===== ROW A -- is it there? decode accuracy, 0 to 1 ====================
  \draw[ink, line width=0.4pt] (0,0) -- (10.5,0);
  \foreach \v in {0,0.25,0.5,0.75,1}{
    \draw[rulegrey, line width=0.3pt] (\xa{\v},0) -- (\xa{\v},-0.10);
    \node[anchor=north, font=\scriptsize, text=quietink] at (\xa{\v},-0.14) {\v};}
  \draw[ink, line width=0.5pt, densely dotted] (\xa{0.083},-0.14) -- (\xa{0.083},0.52);
  \node[anchor=south west, font=\scriptsize\itshape, text=ink]
    at (\xa{0.083},0.52) {shuffled labels};
  % 0.121 -- the constant post-ablation floor. Its label is pushed clear of the
  % shuffled-label rule 0.40cm to its left and leadered back to the mark.
  \fill[quietink] (\xa{0.121},0) circle (1.4pt);
  \draw[quietink, line width=0.3pt] (\xaN{0.121},0.05) -- (1.57,0.20);
  \node[anchor=west, font=\scriptsize, text=quietink]
    at (1.62,0.22) {slab projected out};
  \fill[quietink] (\xa{0.821},0) circle (1.4pt);
  \node[anchor=south east, font=\scriptsize, text=quietink]
    at (\xaL{0.821},0.08) {raw probe, peak};
  \fill[accent] (\xa{0.883},0) circle (1.4pt);
  \node[anchor=south west, font=\scriptsize, text=accent]
    at (\xaP{0.883},0.08) {context removed};
  \node[anchor=east, align=right, font=\scriptsize, text=ink] at (-0.35,0)
    {\textbf{Is it there?}\\ \textcolor{quietink}{decode accuracy}};

  % ===== ROW B -- is it used? ablation cost in multiples of the null =======
  \draw[ink, line width=0.4pt] (0,-2.05) -- (10.5,-2.05);
  \foreach \v in {0,5,10,15,20}{
    \draw[rulegrey, line width=0.3pt] (\xb{\v},-2.05) -- (\xb{\v},-2.15);
    \node[anchor=north, font=\scriptsize, text=quietink] at (\xb{\v},-2.19) {\v};}
  \draw[ink, line width=0.5pt, densely dotted] (\xb{1},-2.19) -- (\xb{1},-1.35);
  \node[anchor=south west, font=\scriptsize\itshape, text=ink]
    at (\xb{1},-1.35) {matched-dimension random subspace};
  % Seven marks, one per layer, and every mark is named. Four are named
  % directly. The other three are layers 16, 12 and 9, reading 3.6, 3.8 and
  % 4.2: 1.1mm and 2.2mm apart on this axis. Three separate leaders into that
  % converge and stop being readable, and fanning them wide enough to separate
  % would claim a resolution the marks do not have, so the cluster is braced
  % and named as an ordered set instead. Reading order is the axis order.
  \foreach \v in {3.6,3.8,4.2,5.7,12.7,14.1,18.6}{
    \fill[accent] (\xb{\v},-2.05) circle (1.4pt);}
  \draw[accent, line width=0.3pt, decorate,
        decoration={brace, amplitude=2.5pt}] (\xb{3.6},-1.93) -- (\xb{4.2},-1.93);
  \node[anchor=south, font=\scriptsize, text=accent, inner sep=1pt]
    at (\xb{3.9},-1.83) {16, 12, 9};
  \foreach \x/\lay in {5.7/6, 12.7/2, 14.1/8, 18.6/4}{
    \node[anchor=south, font=\scriptsize, text=accent, inner sep=1pt]
      at (\xb{\x},-1.83) {\lay};}
  \node[anchor=east, align=right, font=\scriptsize, text=ink] at (-0.35,-2.05)
    {\textbf{Is it used?}\\
     \textcolor{quietink}{ablation cost, $\times$ the null,}\\
     \textcolor{quietink}{marks named by layer}};

  % ===== ROW C -- is it WHICH drug? fraction of the model's own gap ========
  \draw[ink, line width=0.4pt] (0,-5.05) -- (10.5,-5.05);
  % The label is carried separately from the coordinate so the negative tick
  % sets a real minus and not the ASCII hyphen \foreach would print. Same idiom
  % as gate.tex and power.tex; positives stay unmathed so their digits match
  % rows A and B, which are set in text.
  \foreach \v/\t in {-0.02/$-0.02$, 0/0, 0.02/0.02, 0.04/0.04, 0.06/0.06,
                     0.08/0.08, 0.10/0.10}{
    \draw[rulegrey, line width=0.3pt] (\xc{\v},-5.05) -- (\xc{\v},-5.15);
    \node[anchor=north, font=\scriptsize, text=quietink] at (\xc{\v},-5.19) {\t};}
  % The zero rule is drawn in two pieces. The single-cell median sits at
  % -0.0012, one tenth of a centimetre from zero on this scale, so the rule is
  % interrupted across the marker's height rather than run through it.
  \draw[ink, line width=0.5pt, densely dotted] (\xc{0},-5.15) -- (\xc{0},-4.94);
  \draw[ink, line width=0.5pt, densely dotted] (\xc{0},-4.66) -- (\xc{0},-3.34);
  \draw[ink, line width=0.5pt, densely dotted] (\xc{0.10},-5.15) -- (\xc{0.10},-3.34);
  \node[anchor=south east, font=\scriptsize\itshape, text=ink]
    at (\xc{0.10},-3.34) {pre-registered margin};
  \draw[accent, line width=1.0pt] (\xc{0.0063},-3.60) -- (\xc{0.0335},-3.60);
  \fill[accent] (\xc{0.0197},-3.60) circle (1.4pt);
  \node[anchor=west, font=\scriptsize, text=accent] at (\xcR{0.0335},-3.60)
    {residual arm, anchor seed};
  \draw[quietink, line width=1.0pt] (\xc{0.0060},-4.00) -- (\xc{0.0173},-4.00);
  \fill[quietink] (\xc{0.0118},-4.00) circle (1.4pt);
  \node[anchor=west, font=\scriptsize, text=quietink] at (\xcR{0.0173},-4.00)
    {second seed};
  \draw[quietink, line width=1.0pt] (\xc{0.0029},-4.40) -- (\xc{0.0143},-4.40);
  \fill[quietink] (\xc{0.0085},-4.40) circle (1.4pt);
  \node[anchor=west, font=\scriptsize, text=quietink] at (\xcR{0.0143},-4.40)
    {third seed};
  % Solid, like every other point mark; the absent bar is what says it carries
  % no interval. Leader crosses the interruption in the zero rule.
  \fill[quietink] (\xc{-0.0012},-4.80) circle (1.4pt);
  \draw[quietink, line width=0.3pt] (\xcN{-0.0012},-4.80) -- (\xcR{-0.0012},-4.80);
  \node[anchor=west, font=\scriptsize, text=quietink] at (\xcM{-0.0012},-4.80)
    {single-cell arm, median over its depths};
  \node[anchor=east, align=right, font=\scriptsize, text=ink] at (-0.35,-4.20)
    {\textbf{Is it \emph{which} drug?}\\
     \textcolor{quietink}{fraction of the model's}\\
     \textcolor{quietink}{own preference gap}};

  % ===== ROW D -- does a better target help? scramble gap in NIR ==========
  \draw[ink, line width=0.4pt] (0,-7.40) -- (10.5,-7.40);
  \foreach \v/\t in {-0.02/$-0.02$, 0/0, 0.02/0.02, 0.04/0.04}{
    \draw[rulegrey, line width=0.3pt] (\xd{\v},-7.40) -- (\xd{\v},-7.50);
    \node[anchor=north, font=\scriptsize, text=quietink] at (\xd{\v},-7.54) {\t};}
  \draw[ink, line width=0.5pt, densely dotted] (\xd{0},-7.50) -- (\xd{0},-6.12);
  \draw[accent, line width=1.0pt] (\xd{-0.016},-6.35) -- (\xd{0.042},-6.35);
  \fill[accent] (\xd{0.014},-6.35) circle (1.4pt);
  \node[anchor=west, font=\scriptsize, text=accent] at (\xdR{0.042},-6.35)
    {single cell};
  \draw[accent, line width=1.0pt] (\xd{-0.018},-6.75) -- (\xd{0.016},-6.75);
  \fill[accent] (\xd{-0.001},-6.75) circle (1.4pt);
  \node[anchor=west, font=\scriptsize, text=accent] at (\xdR{0.016},-6.75)
    {optimal-transport barycentric};
  \draw[quietink, line width=1.0pt] (\xd{-0.022},-7.15) -- (\xd{0.003},-7.15);
  \fill[quietink] (\xd{-0.010},-7.15) circle (1.4pt);
  \node[anchor=west, font=\scriptsize, text=quietink] at (\xdR{0.003},-7.15)
    {consensus};
  \node[anchor=east, align=right, font=\scriptsize, text=ink] at (-0.35,-6.75)
    {\textbf{Does a better target help?}\\
     \textcolor{quietink}{scramble gap, $\NIR$,}\\
     \textcolor{quietink}{identifiable stratum}};

\end{tikzpicture}

\caption[Four questions, four comparators]{\textbf{Each instrument answers a
different question against a different comparator, and the four rules may not be
read across.} Presence is decode accuracy against a shuffled-label floor; use is
ablation cost in multiples of a matched-dimension random subspace; identity is a
fraction of the model's own clean preference gap, against the $10\%$ margin
pre-registered for calling an effect material; target value is a scramble gap in
$\NIR$ on the identifiable stratum, against zero. Converting one rule into
another is the inference this part exists to block. Bars are $95\%$ intervals
where the instrument returns one; the ablation rule carries none
(\cref{tab:workspace}, \cref{tab:identity}, \cref{tab:epsilon}).}
\label{fig:comparators}
\end{fullwidth}
\end{figure}

% ===========================================================================
\paragraph{The probe, and why its positive is nearly free} The prompt names the
drug in plain text (\cref{sec:methods-data}), so recovering it from the
activations establishes only that the model did not discard what it was handed.
The informative outcome would have been a negative, and what the probe buys is
the depth at which one would appear.\gloss{decode probe}{a classifier trained to
name the drug from the residual stream at one layer}

Forward passes only. Real tier-2 evaluation prompts are grouped by drug; for
each, the residual-stream activation is captured at the last prompt position,
where generation begins, at every layer. A cross-validated multinomial
logistic-regression classifier learns which of the twelve drugs the prompt named,
against the floor measured on shuffled labels ($1/12 \approx 8\%$). The twelve
are a random draw from tier-2 drugs with at least forty scored cells each; that
count sets the $8\%$ floor and, in \cref{sec:res-used}, bounds the drug subspace
at rank $11$. Folds are stratified over prompts and not grouped by well, a
limitation stated and not repaired, and the reason the accuracy is read as a
presence test, not a generalisation estimate.

\paragraph{The negative never arrives} The probe reads $82\%$ at layer 9 and
$76\%$ at layer 16, the deepest measured, against the $\approx 8\%$ floor. The
between-drug to within-drug variance ratio collapses about three-fold across the
final layers, so the drug is present throughout but geometrically de-emphasised
in the layers nearest generation. That de-emphasis is in the raw geometry; once
cell line, plate and dose are projected out, the series rises with depth instead, to
$0.883$ ($88\%$) at the deepest layer. The survival carries the weight, not the
peak, and that is the end of what a probe can say.

\begin{claim}
The drug is not missing from the computation. Its identity is linearly readable
from the residual stream at every layer measured, at the position the first
target token is emitted from. The raw probe reads $82\%$ at its layer-9 peak
and $76\%$ at the deepest against an $8\%$ chance floor; with cell line, plate
and dose projected out it rises with depth instead, to $88\%$
(\cref{sec:res-used}). The prompt names the drug, so this is a floor test and its
positive is close to free. It rules out that the drug field is lost before the
decoder runs, an encoding failure in the activations. It does not rule out that
the decoder ignores it. (\cref{sec:methods-residual} later uses \emph{target
encoding} for a different thing, how the target string is written, and nothing
here bears on that.)
\end{claim}

% ===========================================================================
\beatsection{E}{3}{The drug region is used, at a bit player's weight}{sec:res-used}

\begin{question}[3]
Knowing the drug is written somewhere says nothing about whether the network
reads it. Intervening on the region requires localising it, purifying it of the
batch structure it is confounded with, and proving it really holds the drug
before any null about it can be believed. Generation is sensitive to removing the
region, at a bit player's weight beside cell line.
\end{question}

\noindent
The second instrument is built in four steps, each because the unguarded version
of the same measurement misled us once already.

\paragraph{Step one, where the drug lives} On \num{480} tier-2 prompts (twelve
drugs, forty real cells each) at seven layers, the twelve per-drug mean
activation vectors are centred on the global mean and stacked, and their singular
value decomposition contributes its top ten orthonormal directions,
$V_{\textrm{drug}}$. Twelve centred class means span at most eleven directions,
so the slab keeps essentially the whole span, noisiest included, which for a
removal claim is conservative. The layer grid is even (2, 4, 6, 8, 12, 16) plus
layer 9, added after the raw probe peaked there. A cell-line subspace is built
identically as the positive control.\seealso{the grouping ladder in
\cref{sec:res-blind} is what makes cell line usable as a positive control}

\paragraph{Step two, purification} The raw drug slab is partly a batch slab.
Drugs are confounded with plate and dose by the atlas's design, and probing dose
from inside the drug subspace recovers it at $0.90$--$0.95$ at every layer. The
drug axis is therefore orthogonalised against a context subspace of cell line,
plate and dose directions, giving $V_{\textrm{pure}}$. Every causal claim below
is made on $V_{\textrm{pure}}$, never on the raw slab.

\paragraph{Step three, the removal gate} Before any ``ablating the drug changes
nothing'' result can be believed, the instrument must prove the slab actually
contains the drug. Accuracy is measured before and after projecting the drug
subspace out, and the fraction of above-chance signal removed must exceed $0.8$;
below that the null is unfalsifiable, a measurement of nothing. The gate passes
at every layer. Accuracy falls from $0.41$--$0.82$ to a constant $0.121$ floor,
removing $88$--$95\%$ of the above-chance signal, so the nulls below are real
properties of the model and not broken hooks.

\paragraph{Step four, the causal measurement in logit space} The true sentence is
teacher-forced and its per-token log-probabilities read off clean; a forward hook
then projects the slab out at every position ($h \leftarrow h - (hV)V^{\top}$)
and they are read again. The measure is the mean KL divergence from clean to
ablated over the target tokens, on $60$ prompts.\gloss{teacher forcing}{the
true sentence is scored token by token, never generated} Logit space repairs the
project's first causal attempt, which injected a steering vector and generated;
that vector was small against the residual-stream norm, its effect sat under the
sampling-noise floor, and its ``no effect'' was uninterpretable.\corrmark{the
steering-vector attempt this replaces, \cref{sec:app-identity}} The ablation runs
for $V_{\textrm{pure}}$ (headline), a matched-dimension random subspace (the
null), and the cell-line slab (the positive control). The prompt count is small
and no interval is quoted; what intervals the probe arc has come from the
identity test of \cref{sec:res-mechanism}, which is discounted there in turn, so
no conclusion in this chapter rests on the probe arc alone.

% ===========================================================================
\begin{table}[htbp]
\begin{fullwidth}
\begin{threeparttable}
\caption{\textbf{Drug-label information becomes more decodable with depth while
the same slab carries little functional weight.} All seven layers pass the
removal gate (88--95\% of above-chance drug signal removed); chance is $0.083$.
The last column is ablation cost divided by the matched-dimension random null on
the same layer, a sign-and-magnitude reading and not an estimate with an
interval. The null matches the slab's rank, not the variance it removes --- the
slab's subspace fraction is roughly six times the null's --- so the ratio mixes
drug content with that difference.}
\label{tab:workspace}
\small
\begin{tabular}{@{}lrrrrr@{}}
\toprule
\thead{Layer} & \thead{Decode} & \thead{Decode} & \thead{Ablate pure}
  & \thead{Random} & \thead{Ratio} \\
 & \thead{(raw)} & \thead{(context removed)} & \thead{drug (KL)}
  & \thead{matched dim} & \\
\midrule
2 & 0.408 & 0.354 & 0.155 & 0.012 & $12.7\times$ \\
4 & 0.585 & 0.429 & 1.563 & 0.084 & $18.6\times$ \\
6 & 0.692 & 0.529 & 1.435 & 0.251 & $5.7\times$ \\
8 & 0.779 & 0.569 & 0.596 & 0.042 & $14.1\times$ \\
9 & \textbf{0.821} & 0.602 & 0.418 & 0.100 & $4.2\times$ \\
12 & 0.754 & 0.873 & 0.133 & 0.035 & $3.8\times$ \\
16 & 0.758 & \textbf{0.883} & 0.015 & 0.004 & $3.6\times$ \\
\bottomrule
\end{tabular}
\begin{tablenotes}[flushleft]\footnotesize
\item Decode columns are accuracies over twelve drugs, answering different
questions and not two estimates of one quantity. The raw column includes whatever
cell line, plate and dose contribute; the other is the same probe after those
directions are removed.
\end{tablenotes}
\end{threeparttable}
\end{fullwidth}
\end{table}

% ===========================================================================
\paragraph{What the ablation shows is a dissociation} With cell line, plate and
dose projected out, drug-label decodability rises with depth, from $0.354$ at
layer 2 to $0.883$ at layer 16. Yet the purified slab's descriptive subspace
fraction stays flat at $\approx 0.03$ against $\approx 0.6$ for cell line and
$\approx 0.005$ for the matched-dimension null, an
activation-space summary and not a decomposition of biological drug-response
variance, and ablating it costs $3.6$--$18.6\times$ the matched-dimension random
null. \Cref{fig:mechanism} draws the dissociation whole: the rise, the flat
fraction beneath it, and each layer's ablation cost beside its random null.

That is small beside the cell-line positive control, which the same ablation
reads at $1.0$--$11.2\times$ the drug's. Removing the drug subspace costs
$9$--$37\%$ of what removing cell line costs at layers 2--9, with layer 12 the
single exception at $98\%$ before layer 16 falls back to $25\%$. On subspace
fraction cell line measures $15$--$22\times$ the purified slab. Sensitivity is
established at a bit player's weight, across all seven layers.

\begin{figure}[htbp]
\begin{fullwidth}
\centering
\includegraphics{fig-mechanism}
\caption[Decodable with depth, at a bit player's weight]{\textbf{Drug-label
information becomes more decodable with depth while interventions on the purified
slab have limited effect.} Layer is a true numeric axis in both panels, so the
9--12--16 spacing is drawn as it is.
(a) Top: drug-label probe accuracy over \emph{twelve} drug classes, so chance is
$1/12 = 0.083$; the raw probe is drawn for contrast and the flat grey line is the
same probe after the slab is projected out. Bottom: the descriptive subspace
fraction of the purified drug slab against the cell-line positive control, on its
own separate scale --- it is an activation-space summary, not a decomposition of
biological drug-response variance. The double-headed arrow at layer 16 spans the
two, and is what \emph{encoded, not read} names.
(b) Mean KL from the clean forward pass when the purified drug slab is ablated,
against a matched-dimension random subspace, $n = 60$ prompts per entry. Log
axis because late-layer divergences are small for everything, and paired dots
rather than bars: on a log axis a bar's length is measured from an arbitrary axis
floor, whereas the vertical gap between a pair \emph{is} their ratio, which is
the number annotated beside each pair and the number in
\cref{tab:workspace}. \textbf{No intervals are drawn in either panel.} The
stored dispersions are normal-approximation standard errors over the 60 prompts,
not the thesis's one interval convention (95\% bootstrap clustered over cell
lines); as in \cref{tab:workspace} this is a sign-and-magnitude reading and not
an estimate with an interval. These accuracies come from a different instrument
from the twenty-four-class probe sweep run over the five arms of
\cref{sec:res-mechanism}, whose chance is $1/24$ on a different label set, and
the two must not be read against each other. Whether drug \emph{identity}
is read requires \cref{sec:res-mechanism}.}
\label{fig:mechanism}
\end{fullwidth}
\end{figure}

\paragraph{Where ablation stops} It cannot show that the slab's drug-specific
content is what is read. The purified directions could carry generic
strong-programme energy that merely co-varies with drug across the twelve.
That needs a test which replaces what the slab says.

\begin{claim}
The purified drug-label slab carries decodable information and generation is
sensitive to its removal, at a bit player's weight. Once cell line, plate and
dose are projected out, decodability rises monotonically with depth, which the
raw probe, peaking at layer 9, does not, while the descriptive subspace fraction
stays flat at $\approx 0.03$, roughly $5\%$ of the cell line's $\approx 0.6$, and
is not a biological variance decomposition. Ablating the slab costs $9$--$37\%$
of what ablating cell line costs at layers $2$--$9$, with layer $12$ the
exception at $98\%$. Whether the slab's content is read, ablation cannot say.
\end{claim}

% ===========================================================================
\beatsection{E}{3}{Identity is read, at very low gain}{sec:res-mechanism}
% \claimlede must follow \beatsection: \beatsection clears the stored lede.
\claimlede{The generation decoder is sensitive to the drug-label activation
signal at very low gain, in one training arm and at one rung of the displacement
ladder.}

\begin{question}[3]
Removal cannot distinguish a region whose \emph{content} is read from one whose
mere presence matters. The last question needs an instrument that replaces what
the region says. Push the activations from the drug the prompt names towards
another drug, and ask whether the model moves its preference towards that other
drug's own held-out sentence. It moves, by a sliver, in one arm at one depth. The
design is given at length because three earlier versions could have manufactured
an effect.
\end{question}

\noindent
The verdict first, because the instrument takes several pages to describe and the
verdict is small. In one of the five training arms of \cref{tab:trainconfig}, at
one depth of seven, injecting half the class-mean displacement from drug A to
drug B inside
the purified subspace ($\alpha = 0.5$, the rung the multiplicity correction is
applied to) moves the model's preference towards drug B's own held-out sentence,
against an identically constructed injection naming a third drug, by
\ci{+0.0197}{+0.0063}{+0.0335} of the preference gap the model already has. One
to two per cent, against a $10\%$ margin pre-registered for calling an effect
material. Two further seeds reproduce it, and of twenty-six headline tests it is
the only survivor of multiplicity correction.

Read at face value that says the generation decoder does consult which drug the
region names, at a gain far too small to account for anything in
\cref{sec:res-blind}. Read with the multiplicity discipline below it says less.

\begin{scopelimit}[carried into \cref{sec:res-epsilon} and \cref{sec:limitations}]
This is the weakest load-bearing measurement in the chapter and it is carried as
such. No argument downstream rests on it alone, and a reader who discards it
entirely loses nothing that the drug-blindness result depends on. What the
following section tests is a prediction of the \emph{interpretation} of this
result, not the result itself, and the inference there is written to run one way
only for that reason.
\end{scopelimit}

\paragraph{First, a contrast score} It is
\begin{equation}
 \log P\!\left(g_B \mid \textrm{prompt}_A\right) \;-\;
 \log P\!\left(g_A \mid \textrm{prompt}_A\right),
\end{equation}
where $g_A$ and $g_B$ are real cell sentences emitted under drugs A and B by
held-out cells, each log-probability a per-token mean. Any offset uniform across
positions, an inflated logsumexp normaliser included, then cancels exactly. It
does \emph{not} cancel in the stronger sense one might assume. Once $g_A$ and
$g_B$ diverge the two teacher-forced sequences visit different states, so the
normaliser's state-dependent part survives, unbounded here. The \texttt{permuted}
arm controls for that, sharing anchor, construction and norm and differing only
in naming the wrong drug. What survives the two-arm comparison is the model's
preference between drug B's own held-out sentence and drug A's.

\paragraph{Second, one anchor and one norm} Writing $\Pi$ for the orthogonal
projector onto the purified slab, the \texttt{swap} arm injects
$\Pi(\mu_B - \mu_A)$, pointing at the drug whose sentence is scored; the
\texttt{permuted} arm injects $\Pi(\mu_C - \mu_A)$ at identical norm, same anchor
and construction, wrong drug. An isotropic random direction is kept only as a
damage diagnostic, since it can disrupt the model in ways an observed class-mean
displacement does not. The subtraction cancels the global mean exactly, and
nothing is ablated, so the control has no removed term to trip over.

\paragraph{The three failed designs taught the constraints} A control drawn in
the full $2048$-dimensional stream puts $\approx 99.5\%$ of its energy where the
intervention never acts, so it must live in the treatment's subspace.
Ablate-then-add delivers (added $-$ removed), so an uncentred class mean turns
substitution into restoration, hence no ablation and an exactly centred
displacement. A norm-matched random direction inflates the logsumexp normaliser
(measured $+2.75$ nats), depressing every target whether or not it carries
identity. The planted-world selftest validates the directional logic but cannot
speak to the normaliser, since there each drug's sentence is a single token.
Algebra in \cref{sec:app-identity}.

\paragraph{Out of sample, and on a comparable scale} Subspaces and class means are
fit on one slice of cells; $g_B$ comes from a third disjoint slice, never from
the rows defining $\mu_B$, else the injected vector and the scored sentence would
share cells and the test would be partly circular. Because the five arms emit
different target formats, every effect is divided by that model's own clean
preference gap, measured with no intervention.

\paragraph{Four guards carry the inference} A drug-proxy screen and an
injected-norm diagnostic are engineering, deferred to \cref{sec:app-identity}.

The first is the uncertainty model, a cluster bootstrap over ordered (A,B) pairs,
because rows sharing a pair share the injected vector and the scored sentence. At
the anchor seed's layer $12$ that is $60$ rows in $58$ pair clusters over $24$
anchor drugs, close to a row bootstrap. It guards against sharing within a pair,
not across pairs sharing an anchor or target drug.\gloss{cluster
bootstrap}{resamples whole (A,B) pairs, not rows, so a shared vector
counts once}

The second is the equivalence criterion, stricter than TOST. The interval must
lie inside a pre-registered margin, set at $10\%$ of the model's own clean
preference gap, \emph{and} contain zero. That conjunct runs one way only; it can
withhold an equivalence verdict, never manufacture one. Under textbook TOST the
headline \ci{+0.0197}{+0.0063}{+0.0335} lies wholly inside the margin and would
be returned as equivalent, read as a null; this rule returns a directional effect
smaller than the margin.

The third is a per-layer headroom gate, excluding any layer where the cell-line
positive control shows less than $5\times$ dynamic range, because a saturated
instrument can support neither verdict. The fourth is a magnitude ladder,
repeating the test at $\alpha \in \{0.5, 1, 2, 5, 10\}$ times the natural
displacement, because a real identity effect should show a monotone
dose--response. The ladder is the guard that bites.

\paragraph{One survivor in twenty-six tests} The design ran across the five
training arms of \cref{tab:trainconfig} and, for the two residual arms, at three
further seeds. Of $26$ headline tests, five arms by up to seven depths, exactly
one survives Benjamini--Hochberg correction, the residual arm at layer $12$
($q = 0.026$, pooling the $26$ across arms and not correcting within each as the
pipeline stores them). That $q$ rests on a bootstrap $p$, and the arm's probe
file carries two draws of it for the same effect, identical to sixteen digits, at
$p = 0.001$ and $p = 0.004$, either side of the $0.05/26$ admitted at rank one;
on the second draw nothing survives. The consensus arm at layer $2$ ($p = 0.012$)
does not ($q = 0.156$).

\paragraph{Which rung the correction is applied to matters} The headline is the
$\alpha = 0.5$ rung, half the natural displacement, so the ladder's first entry
is that headline in raw nats ($+0.0013$) and not an independent point. At full
displacement the effect is \emph{larger} and noisier,
\ci{+0.0265}{+0.0021}{+0.0524} against \ci{+0.0197}{+0.0063}{+0.0335}, at
$p = 0.033$ against $p = 0.001$ on the draw the correction uses. Each of
$\alpha = 0.5$, $1$ and $2$ excludes zero; $\alpha = 5$ and $10$ span it. In raw
nats the rungs read $+0.0013$, $+0.0017$, $+0.0044$, $+0.0061$, $+0.0046$, rising
to $\alpha = 5$ and turning over at $\alpha = 10$ (\cref{fig:family}, lower
strip), consistent with a dose--response without establishing one. All five
re-score the same $60$ rows at scaled injection, no trend statistic is computed,
and the pre-registered monotonicity holds over four rungs of five.

But applying the $26$-test correction to the $\alpha = 1$ slice leaves
\emph{nothing} surviving, and applying it to $\alpha = 2$ promotes a different
arm at a different depth (\texttt{residual\_holdout} at layer $4$, $q = 0.007$)
while demoting this one (\cref{fig:family}). The slice decides not merely
whether a survivor exists but which one, and below that the bootstrap draw
decides whether one exists at all. That is the mark of a family dominated by
noise, not of one containing a single robust effect. The multiplicity-corrected
claim is therefore a statement about the $\alpha = 0.5$ slice and is not robust
to that choice, which is why this section opened by discounting its own result
instead of closing on it.

% ===========================================================================
\begin{figure}[htbp]
\begin{fullwidth}
\centering
\includegraphics{family}
\caption[The single survivor moves with the rung]{\textbf{The one test that
survives multiplicity correction changes arm and depth with the rung of the
displacement ladder, and at one rung there is no survivor at all.} Each panel is
one slice of the ladder: the five training arms of \cref{tab:trainconfig} (rows)
by the seven probed depths (columns), $n = 26$ measured tests per slice, on one
diverging scale shared across the three panels. Colour is the \texttt{swap}
minus \texttt{permuted} effect as a fraction of that arm's own clean preference
gap. The layer axis is ordinal, drawn at equal spacing, and is not proportional
to depth. Nine of the $35$ grid positions were never measured and are drawn as
empty dotted outlines, not as zeros. Benjamini--Hochberg is pooled across arms
within a slice, as the text applies it and not as the pipeline stores it; $q$ is
printed in the three cells that fall below $0.10$ and ringed where $q < 0.05$.
That criterion admits the residual arm at layer $12$ at $\alpha = 0.5$
($q = 0.026$, on the stored bootstrap draw the text discusses), nothing at
$\alpha = 1$, and \texttt{residual\_holdout} at layer $4$ at $\alpha = 2$, which
demotes layer $12$ to $q = 0.058$; it is the only survivor among this chapter's
twenty-six headline tests. Beneath, the same residual arm at layer $12$ in raw
nats across all five rungs, with the $10\%$ pre-registered materiality margin as
the vertical rule. Intervals throughout are $95\%$ cluster bootstraps over
ordered (A,B) pairs, $4000$ resamples, which at this layer is $60$ rows in $58$
pair clusters over $24$ anchor drugs. The three seed replicates are in
\cref{tab:identity} and are deliberately not drawn here.}
\label{fig:family}
\end{fullwidth}
\end{figure}

% ===========================================================================
\begin{table}[htbp]
\begin{fullwidth}
\begin{threeparttable}
\caption{\textbf{The identity test replicates in the arm that behaviour calls
drug-sensitive and is absent in the arm behaviour calls drug-blind.} Effects are
a fraction of that model's own clean preference gap between the two sentences
compared (mean effect $+0.0133$ across the three seeds that measured it). The $p$
column is a two-sided cluster bootstrap over $4000$ resamples floored at
$1/4000$; seed 43's entry is that floor, reached when no resample fell at or
below zero, so it is a resolution limit and not a measurement. Seed 44 produced
no layer-12 measurement, and five of its seven layers were skipped because the
purified slab retained context, not for want of headroom; those five carry
headroom of $9.7$ to $23.3$. The pre-registered criterion of two replications in
three new seeds is met, the third a non-measurement and not a failure.}
\label{tab:identity}
\small
\begin{tabular}{@{}lrrr@{}}
\toprule
\thead{Arm / seed} & \thead{effect} & \thead{95\% CI} & \thead{$p$} \\
\midrule
residual, seed 42 & $+0.0197$ & $[+0.0063, +0.0335]$ & $0.0010$ \\
residual, seed 43 & $+0.0118$ & $[+0.0060, +0.0173]$ & $0.00025$ \\
residual, seed 45 & $+0.0085$ & $[+0.0029, +0.0143]$ & $0.0015$ \\
\midrule
single-cell arm, six depths
  & \multicolumn{3}{l}{nothing detectable at any depth (median $-0.0012$)} \\
\bottomrule
\end{tabular}
\begin{tablenotes}[flushleft]\footnotesize
\item The rows are the headline $\alpha = 0.5$ rung, to which the multiplicity
correction that admits the first row is applied and to no other.
\end{tablenotes}
\end{threeparttable}
\end{fullwidth}
\end{table}

% ===========================================================================
\paragraph{The instrument has been shown to register a positive} The residual arm
is independently known to respond to the drug its prompt names. On \texttt{train},
swapping the named drug and its mechanism for an orthogonal partner costs it
$+0.0898$ in $\NIR$, excluding zero clustered on cell line
($[+0.0382, +0.1415]$) and on drug ($[+0.0291, +0.1506]$), a behavioural result
established in \cref{sec:res-generalisation}. \forward{sec:res-generalisation}

Three things bound what it licenses. It is a training-split fact, and does not
extend to held-out combinations. It is behavioural, so it says the arm responds
to the drug its prompt names, not that the response is routed through the
subspace this test injects into. And the two measurements are not on the same
weights, the swap gap scored on the rebuilt target build and this probe on the
earlier one (\cref{tab:trainconfig}), so drug-sensitivity and identity reading
are attributed to residual-target training and not to any checkpoint. Narrower
than a routing claim, and enough. The probe detects identity reading, replicably,
in the arm behaviour calls drug-sensitive, and none at any depth in the arm that
is behaviourally drug-blind. A null from an instrument never shown to register a
positive is uninterpretable; this one has.

\paragraph{Three limits belong with the verdict} The effect appears at one depth
of seven, and layer 4 is positive in a single seed and null in the other three.
The single-cell arm has no layer-12 measurement, its cell-line positive control
reaching $3.88\times$ there, under the $5\times$ gate; so the cleanest
comparison, the same depth across two arms, does not exist, and the dissociation
rests on that arm being null at all six depths it could measure. And the sibling
residual-holdout arm does not reproduce the effect at the headline rung
($+0.0028$ and $+0.0049$ in the two seeds where it was measured); it surfaces
instead at $\alpha = 2$, rung-dependence and not a second replication.

Read with the decodability series, the picture is label information that becomes
more linearly decodable with depth once cell line, plate and dose are projected
out, while the process that writes the prediction barely consults it. That is
suggestive of the workspace account of \cref{sec:lit-interp}, an interpretation
of the pattern and not a demonstration of the mechanism.

\begin{claim}
The anchor seed puts it at \ci{+0.0197}{+0.0063}{+0.0335} of the model's own
preference gap, and $+0.0133$ averaged over the three seeds that measured it, one
to two per cent, a fifth or less of the $10\%$ margin pre-registered for calling
an effect material. The multiplicity-corrected version of that holds for the
$\alpha = 0.5$ rung on one of its two stored bootstrap draws, not for its
neighbours. At $\alpha = 1$ nothing survives correction; at $\alpha = 2$ a
different arm does. The claim is not that the decoder cannot read drug identity,
and it is not the stronger mechanistic story earlier drafts told.
\end{claim}

% ===========================================================================
\beatsection{E}{4}{Denoising the target changes the target, not the decoder}{sec:res-epsilon}

\begin{question}[4]
If the defect is a decoder that barely consults the drug region whatever is
written there, then improving what sits in that region should not help, and
improving the training target is exactly that kind of intervention. It is also
the repair a reader reaches for first, which is reason enough to test it. Three
target constructions, one scoring, one stratum. None separates from a prompt
naming the wrong drug.
\end{question}

\noindent
The interpretation at the end of the last section is \notestablished, and this
section does not proceed as though it were. The inference runs one way only. A
positive here would tell against the low-gain reading; a negative is evidence
about target-side fixes, not confirmation of the mechanism. Because one failed
construction proves nothing, two more were compared with it.

\paragraph{Three points, not a sweep} They are the original single-cell target,
one optimal-transport barycentre and the consensus target. \emph{Consensus}
groups cells by (drug, cell line, dose) and averages them into one pseudobulk
sentence, every original prompt retained so only the target changes.
\emph{Optimal transport} uses the full condition key carried by the source,
(cell line, plate, drug, dose/treatment well), computing each condition
separately, its treated cells coupled only to the shared DMSO controls from that
condition's own (cell line, plate). Treatment conditions are never pooled.

The coupling $\pi$ minimises $\sum_{ij}\pi_{ij}\lVert x_i - y_j\rVert^2$ under
balanced uniform marginals, which impose mass conservation between the empirical
control and treated measures, a modelling assumption and not an observed
biological constraint. Sinkhorn iteration~\cite{cuturi} solves it on the kernel
$\Gamma = \exp(-C/\varepsilon)$, with $C$ the pairwise squared-distance matrix.
Couplings are computed in the released Tahoe scVI latent~\cite{scvi}
($n_{\textrm{latent}} = 10$), joined per cell by barcode and not re-encoded,
since a re-encode produced a degraded latent (cell-line probe $0.11$ versus
$0.975$). Targets are built without the scVI decoder, as convex combinations of that
condition's observed treated cells, so they stay inside their convex hull, not
necessarily on the data manifold.

The regularisation parameter locates the middle construction between two limits.
As $\varepsilon\rightarrow\infty$ the coupling is uniform and the target is the
treated mean (consensus); as $\varepsilon\rightarrow 0$ it approaches an
unregularised sparse transport solution, not necessarily a map to one observed
treated cell. Only three hand-picked constructions were trained and scored, so
this is a three-point comparison and not a sweep of the family.

\paragraph{One scoring, named before the result} Every point was scored on the
identifiable stratum, the conditions \cref{sec:res-blind} selected because their
own within-well reference clears the pre-set $0.80$ bar, under the random-partner
scramble defined there. That is the stratum on which a zero-drug-information
control-copy already matches the model, so the scope limit of
\cref{sec:methods-probe} applies to every row of \cref{tab:epsilon}.

% ===========================================================================
\begin{table}[htbp]
\begin{fullwidth}
\begin{threeparttable}
\caption{\textbf{Three target constructions, and no point separates from its
random-partner scramble on the identifiable stratum} ($n = \num{204}$). Two
things bound that negative. The consensus endpoint was stopped at roughly $60\%$
of an epoch with its tier-1 NIR unscoreable, so the negative rests on the
single-cell and optimal-transport endpoints. And the optimal-transport endpoint
is not null in every frame. These rows use a random partner at $1200$ generation
tokens, while against the geometry-defined \texttt{orth} partner over all its
conditions the same checkpoint moves from \ci{+0.0144}{-0.0051}{+0.0338} at $600$
($n = \num{250}$) to \ci{+0.0292}{+0.0144}{+0.0441} at $1400$ ($n = \num{248}$),
reported in \cref{sec:res-encoding}. Comparator, population and budget differ at
once, and \cref{sec:methods-data} shows the budget is a claim-level variable, so
the frames are set side by side and not put into ratio. This is a null under this
scoring, not over the family of target constructions.}
\label{tab:epsilon}
\small
\begin{tabular}{@{}llll@{}}
\toprule
\thead{Target} & \thead{$\varepsilon$} & \thead{$\gap$} & \thead{Verdict} \\
\midrule
single cell (real treated cell) & $\approx 0$
  & \ci{+0.014}{-0.016}{+0.042} & not separated \\
OT barycentric & $0.5$
  & \ci{-0.001}{-0.018}{+0.016} & not separated \\
consensus (global mean) & $\rightarrow\infty$
  & \ci{-0.010}{-0.022}{+0.003} & not separated \\
\bottomrule
\end{tabular}
\begin{tablenotes}[flushleft]\footnotesize
\item $\gap$ is the model's $\NIR$ minus its scramble's, so its null is zero and
not chance. The three rows share a scoring, a stratum and a comparator; they are
one comparison, not three independent experiments. All three intervals are
\emph{one-way} percentile cluster bootstraps over the $40$ cell lines, not the
two-way default.
\end{tablenotes}
\end{threeparttable}
\end{fullwidth}
\end{table}

% ===========================================================================
An interval that covers zero earns its keep by what it excludes. On the
identifiable stratum the single-cell arm's scramble gap is
\ci{+0.014}{-0.016}{+0.042}, against the $0.170$ separating the model's $0.768$
from that stratum's $0.938$ reference. At the very top of the interval, naming
the right drug instead of a wrong one would buy a quarter of what is there to be
bought, and the point estimate puts it at $8\%$ of that. Taking that distance at
face value is the generous reading, for the reason declared in
\cref{sec:methods-probe}. The interval does not show the effect is zero and no
experiment here could; it puts it far below the room these conditions leave.

% ===========================================================================
\paragraph{Training was not inert, and what it did instead is a result} The
consensus model stops echoing the control. On the identifiable stratum it scores
$0.653$ against a zero-drug-information control-copy at $0.766$, where the
single-cell model sits on the control at $0.768$. Its predictions are not
mode-collapsed either; there is real between-condition variation. But that
variation is uncorrelated with real drug identity. The Mantel correlation between
the model's between-drug distance structure and the true one averages $+0.03$
over cell-line and plate groups.

The failure this literature names first for a perturbation model that
underperforms on well-calibrated metrics is mode collapse, where the prediction
ignores the input and reverts to a mean, and the repair is to restore diversity
in what the model emits~\cite{mejia}. This arm is not that. It moved off the control, kept the
variation, and spent it on the wrong axis. Producing no structure and producing
structure that does not track the drug are different failures with different
repairs.

\begin{claim}
No point in the three-point comparison separates from its scramble on the
identifiable stratum under the random-partner test. Two things bound that
negative. The consensus arm was stopped mid-epoch, so the comparison rests on the
single-cell and optimal-transport points, and the optimal-transport arm is not
null in every frame: scored against the geometry-defined \texttt{orth} partner
over all its conditions at the residual arm's generation budget, its gap excludes
zero (\cref{sec:res-encoding}). By the inference rule above that positive tells
against the low-gain reading and not for it, and being in a different frame
it does not settle the matter, but it is one more reason the interpretation of
\cref{sec:res-mechanism} is carried here as an interpretation and not a result.
What the comparison establishes is that denoising the target changes what is in
the target without changing what the decoder does with it under this scoring, and
not that no target-side fix can work. Measuring how much of the target its tokens
expose to the drug, which the next two sections do, is not the same as showing
that the target encoding is the defect (\cref{sec:methods-residual}). That
construction changes five things at once.
\end{claim}
